\chapter{Dysgeusia}

\section*{Definition}
Dysgeusia is a distortion of the sense of taste, most commonly described as a persistent metallic, bitter, salty, or rancid flavor in the mouth that is independent of food intake.

\section*{When to See a Doctor}
See your doctor if:
\begin{itemize}
  \item Dysgeusia persists for \textgreater{}2--4 weeks after the causative medication or infection is resolved
  \item Associated with unintentional weight loss, reduced appetite, or malnutrition
  \item Accompanied by focal neurologic deficits, facial weakness, or loss of smell (anosmia)
  \item Oral pain, burning, or visible lesions are present
  \item Onset is sudden or temporally linked to a new medication
\end{itemize}

\section*{How It Develops}
Changes in taste may result from altered taste receptors, smell, saliva, oral disease, medicines, infection, or injury to sensory pathways. Because smell contributes substantially to flavor, a smell disorder may be perceived as a taste problem.

\section*{Causes}
\begin{itemize}
  \item \textbf{Pharmacologic:} Chemotherapy (cisplatin, carboplatin, 5-FU), ACE inhibitors (captopril), metronidazole, clarithromycin, lithium, terbinafine, topiramate, methimazole
  \item \textbf{Oral/dental:} Xerostomia (Sj\"{o}gren syndrome, radiation therapy, medications), periodontal disease, oral candidiasis, lichen planus, geographic tongue
  \item \textbf{Neurologic:} Bell palsy, multiple sclerosis, stroke (especially thalamic or insular), Guillain--Barr\'{e} syndrome, Parkinson disease
  \item \textbf{Nutritional:} Zinc deficiency (most common), copper deficiency, vitamin B\(\_{12}\) deficiency
  \item \textbf{Metabolic:} Diabetes mellitus, chronic kidney disease, hepatic failure, hypothyroidism, Cushing syndrome
  \item \textbf{Infectious:} Upper respiratory infections, SARS-CoV-2 (COVID-19), influenza, sinusitis, oral HSV
  \item \textbf{Miscellaneous:} Aging (presbygeusia), smoking, radiation therapy to head/neck, idiopathic
\end{itemize}

\section*{Related Symptoms}
\begin{itemize}
  \item Anosmia or hyposmia, xerostomia, oral burning (burning mouth syndrome), halitosis, dysphagia, weight loss, nausea, food aversions, depression, decreased appetite
\end{itemize}
\textbf{Important:} Because taste and smell are intimately linked (retronasal olfaction), many cases of perceived dysgeusia actually reflect olfactory dysfunction; a thorough evaluation of both chemosensory systems is warranted.

\section*{Medical Evaluation}
\begin{itemize}
  \item Detailed medication review and timeline; trial discontinuation of offending agent if safe
  \item Formal taste testing is available in specialist settings but is not routine for every patient
  \item Targeted laboratory studies are guided by the history and examination; routine broad testing for zinc, copper, and multiple vitamins is not required for everyone
  \item Oral examination: inspect oral mucosa, tongue, dentition; check for candidiasis (may be subtle)
  \item MRI brain with thin cuts through the olfactory bulbs and insula if neurologic symptoms present
\end{itemize}

\section*{Treatment Options}
\begin{itemize}
  \item \textbf{Nutritional:} Replace a documented nutrient deficiency under professional guidance; prolonged high-dose zinc can cause copper deficiency
  \item \textbf{Oral care:} Treat candidiasis with topical antifungals (clotrimazole troches, nystatin); artificial saliva products for xerostomia; pilocarpine or cevimeline for severe dry mouth
  \item \textbf{Medication management:} Substitute offending drug when possible (e.g., switch ACE-I to ARB, metronidazole to alternative antibiotic)
  \item \textbf{Dietary counseling:} Nutritionist referral; high-protein supplements if intake is reduced; avoid extreme food temperatures that may worsen perception
  \item \textbf{Neurologic:} Corticosteroids for Bell palsy or inflammatory causes; gabapentin or alpha-lipoic acid for burning mouth syndrome
\end{itemize}

\section*{Self-Care and Home Management}
\begin{itemize}
  \item Practice meticulous oral hygiene: brush twice daily, floss, and scrape the tongue
  \item Use sugar-free gum or mints to stimulate saliva and mask metallic taste
  \item Rinse the mouth with a baking soda solution (1/4 tsp in 8 oz water) before meals
  \item Experiment with flavor enhancers: lemon juice, vinegar, herbs (basil, oregano), soy sauce, or strong spices (cinnamon, mint)
  \item Avoid smoking, alcohol-based mouthwashes, and metallic cookware (aluminum, copper)
  \item Stay well hydrated to maintain salivary flow
\end{itemize}

\section*{References}
\begin{thebibliography}{99}
\bibitem{dysgeusia:ref1} Bromley SM. ``Smell and Taste Disorders: A Primary Care Approach.'' Am Fam Physician. 2000;61(2):427--436.
\bibitem{dysgeusia:ref2} Malaty J, Malaty IAC. ``Smell and Taste Disorders in Primary Care.'' Am Fam Physician. 2013;88(12):852--859.
\bibitem{dysgeusia:ref3} Doty RL, Bromley SM. ``Taste and Smell Disorders.'' In: Aminoff MJ, Boller F, Swaab DF, eds. Handbook of Clinical Neurology. Elsevier; 2008.
\bibitem{dysgeusia:ref4} Deems DA, Doty RL, Settle RG, et al. ``Smell and Taste Disorders: A Study of 750 Patients.'' Arch Otolaryngol Head Neck Surg. 1991;117(5):519--528.
\bibitem{dysgeusia:ref5} Henkin RI, Schecter PJ, Friedewald WT, et al. ``A Double Blind Study of the Effects of Zinc Sulfate on Taste and Smell Dysfunction.'' Am J Med Sci. 1976;272(3):285--299.
\end{thebibliography}
